% discussion.tex (국문 v15, 2026-10-21): 영문 v15 의 쌍둥이. 소제목 없이 여섯 문단: 기준값의 역할, 알래스카 이득 부재의 이유, 제품 비교의 의미, 선행 연구와의 관계, 지도와 새 지역, 한계와 전망.
\section{고찰}\label{sec:discussion}

주요 결과는 물리 정보를 모델의 어느 위치에 넣는지에 따라 지역 안 정확도와 지역 간 전이성이 달라진다는 것이다. 물리 잔차 구성(식~(\ref{eq:residual}))에서 Stefan 항은 각 연도의 융해 도일에 대한 ALT 의 반응을 TDD 의 제곱근에 비례하는 형태로 나타내고, 학습기는 토양과 지형 때문에 지역 계수 하나로 설명되지 않는 부분을 보정한다. 이 역할 분담이 결과의 주요 양상을 설명한다. Stefan 모델 대비 이득은 지형과 토양 공변량이 있을 때만 나타났으며, 이는 이 변수들이 지점별 유효 계수를 좌우한다는 해석과 부합한다. 이득은 Stefan 모델의 오차가 큰 조건일수록 컸다. 잔차 모델의 SHAP 기여도에서 지형 변수는 레나 델타에서 \nshapLeResTerrain\%, 토양 변수는 캐나다에서 \nshapCaResSoil\% 를 차지하였다. 세 양상은 잔차 학습기가 Stefan 계수 하나로 설명되지 않는 융해 깊이의 공간 변동을 학습하였음을 시사한다. 기후 응답을 라벨에서 배워야 하는 직접 ML 은 한 지역의 라벨이 수천 개일 때 지점-연도 단위 RMSE 가 더 낮았으나, 라벨 25개로는 알래스카와 캐나다에서 Stefan 모델보다 부정확하였고 레나 델타를 학습에서 빼면 그 오차가 Stefan 모델보다 \ndrefLeregDzabs{} cm 컸다. 물리 잔차 모델의 지역 안 RMSE 는 직접 ML 보다 최대 \ndrefRoDzblkmax{} cm 컸으나 기준값의 전이성을 유지하였다. 이 성질은 라벨이 적은 지역에서 중요하다.

알래스카에서는 잔차 보정에 의한 이득이 없었으며, 그 원인은 두 가지로 해석된다. 알래스카에서는 재보정 Stefan 모델이 이미 지역 평균 수준과 연도 간 변동을 재현하므로 공변량으로 설명할 체계적 잔차가 거의 남지 않는다. 또한 남은 오차의 대부분은 한 1 km 셀 안 탐침 지점 사이의 산포(합동 표준편차 \ncyWithinSdAk{} cm)이며, 1 km 해상도의 어떤 공변량도 이를 해소하지 못한다. 모든 삼분위 구간에서의 이득 부재(그림~\ref{fig:conditions}b)와 셀-연도 평균에서의 같은 결과(보충 표~S5)는 이 두 원인과 부합하며 학습기의 결함을 뜻하지 않는다. 따라서 실제 적용에서 학습기가 늘 필요한 것은 아니다. 공간 블록 검증에서 물리 잔차 모델과 재보정 Stefan 모델을 비교하면 그 지역에 공변량으로 설명되는 체계적 잔차가 있는지(레나 델타 · 캐나다처럼 판정된 음의 $\Delta$RMSE) 없는지(알래스카)를 확인할 수 있으며, 잔차 학습기는 앞의 경우에 적용한다.

같은 지점 · 같은 해에서 비교하면 ML 제품의 정확도는 학습 지점과의 거리에 따라 달랐으며, 이 차이는 제품 자체의 지점 집합 안 검증으로는 평가되지 않는다. 제품의 학습 지점이 우리 라벨 지점 가까이 있는 곳(알래스카, 그리고 같은 현장 관측의 이전 판을 학습 자료에 포함한 Liu 등의 경우 캐나다)에서 ML 제품은 우리 모델만큼 정확하였고, 그곳의 비교는 독립적이지 않다. 근처에 학습 지점이 없는 곳(레나 델타, 중앙값 거리 \ndistweiLemed{} km)에서 두 ML 제품은 관측 ALT 와 무관하게 약 40--70 cm 를 예측하였고, 레나 델타 전체에서 약 20 cm 과대 추정하였으며(델타 북부의 얇은 활동층에서 가장 큼), 레나 델타 라벨 없이 학습한 물리 잔차 모델보다 부정확하였다. 이 양상은 지역 홀드아웃에서 직접 ML 이 보인 오차 증가와 같은 형태이며, 두 ML 제품이 연도별 예측 변수를 쓰지만 기후와 ALT 의 관계를 다른 지역의 학습 지점에서 학습하였기 때문으로 해석된다. 물리 잔차 모델에서는 명시된 해의 TDD 에 대한 반응을 학습 지점과 무관한 Stefan 항이 나타내므로 이러한 오차 증가가 작다. ALT 관측으로 학습하지 않는 과정 기반 ESA CCI 제품은 세 지역 모두에서 우리 모델보다 부정확하였으나, ML 제품과 달리 레나 델타에서 체계적으로 치우치지 않았고(그림~\ref{fig:products}c), 이는 그 물리적 기반과 부합한다.

이 결과는 선행 비교를 확장한다. Gautam 등\cite{Gautam2025}은 한 해의 알래스카 CALM 지점 68개에서 랜덤 포레스트가 홀드아웃 지점에서 정확도의 대부분을 잃는 반면 Stefan 모델은 이를 유지함을 보였다. 이 순서는 세 지역의 2010--2024년 지점-연도와 지역 홀드아웃에서도 유지되며, Stefan 모델을 보정하는 학습기가 전이성을 지키므로 두 접근은 상보적이다. 호수와 하천의 물리 기반 학습 연구는 과정 지식이 주로 자료 부족과 외삽에서 도움이 된다고 보고하였고\cite{Read2019,Jia2021,Willard2022}, 이 연구의 라벨 수 실험과 지역 홀드아웃 실험은 이를 ALT 에 대해 정량화한다. 무작위 교차검증이 직접 ML 의 공간 블록 오차를 \noptimLeDz--\noptimCaDz{} cm 과소평가한다는 결과는 공간 검증 문헌\cite{Roberts2017,Ploton2020,MeyerPebesma2022}과 일치한다. 공간 블록 검증과 지역 홀드아웃 결과를 함께 보면, 무작위 또는 지점 기반 분할로 얻은 오차는 학습 관측이 없는 지역에서의 제품 정확도를 나타내지 못한다.

연도별 지도는 물리 잔차 모델을 세 지역에 적용한 산출물이며, 적용 가능 영역과 예측 구간으로 사용 조건을 명시한다. 각 지도는 해당 연도의 ERA5-Land 기후(9월 1일까지 누적한 TDD)를 입력으로 하므로 같은 해의 관측이나 제품과 비교할 수 있다. 지도의 지역 평균 편차는 융해 계절 기후에 대한 모델 응답이다(지역 TDD 와의 상관 \ntddanomAk{}--\ntddanomCa). 이 편차는 해마다 측정되는 지점에서만 관측과 대조할 수 있다(보충 그림~S9). 지도의 선형 추세는 2011--2024년 기후에 대한 모델 응답이며 관측된 추세가 아니다. 지역 관측이 없는 경우에는 다른 지역의 계수를 쓴 Stefan 모델이나 다른 지역으로 학습한 물리 잔차 모델을 출발점으로 쓸 수 있다. 두 모델은 레나 델타에서 직접 ML 과 같은 큰 오차 증가를 보이지 않았으나, 캐나다에서는 다른 지역으로 학습한 직접 ML 이 가장 정확하였다. 지역 라벨 25개의 효과는 지역에 따라 달랐다. 물리 잔차 모델의 오차는 레나 델타에서 \ntlrmseLeRszero{} cm 에서 \ntlrmseLeRstwfive{} cm 로 줄었으나, 캐나다에서는 \ntlrmseCaRszero{} cm 에서 \ntlrmseCaRstwfive{} cm 로 늘었다. 캐나다의 증가는 두 자료 출처 사이의 융해 깊이 차이를 라벨 25개로 맞춘 계수 하나가 나타내지 못하기 때문인 것으로 추정된다. 두 경우는 새 지역 안의 공간 블록 비교와 결과 지도의 적용 가능 영역으로 구별할 수 있다.

이 연구에는 한계가 있다. 알래스카 라벨의 절반이 2014년에서 나와 그 지역에서는 연도별 입력의 효과를 판정하지 못하였고, 레나 델타 라벨은 \nnblkLe{}개 블록의 한 자료에서 나와 구간이 넓다. 라벨은 세 프로그램의 탐침 측선과 CALM 격자 평균을 합친 것이며 모델 오차는 프로그램에 따라 다르다(보충 표~S7). 특히 캐나다 CALM 지점-연도 \nsrcCaCALMN{}개는 평균 ALT \nsrcCaCALMObs{} cm 가 \nsrcCaCALMBias{} cm 과소 추정되어, 측정 규약과 지점 선택의 차이가 보고된 오차에 포함된다. 라벨은 독립적인 1~km 셀이 아니라 관측 지점이며, 공간 블록 설계와 셀-연도 평균 평가가 이를 고려한다. 기후 입력은 월 단위이고 토양 공변량은 정적이다. 눈과 수분 항을 포함하는 Kudryavtsev 모델은 캐나다에서 재보정 Stefan 모델보다 정확하였고, 그 차이의 대부분은 잔차 학습기가 공변량으로 보정하였다. 수분을 반영한 기준값은 잔차 보정 뒤에 남는 캐나다의 오차를 더 줄일 가능성이 있다. 약 3~m 보다 두꺼운 활동층은 학습 범위 밖이고, 한 해의 지도와 관측 비교는 레나 델타와 캐나다에서 소수의 셀로 이루어진다. 이 연구의 모든 수치가 명시된 지점 · 해 · 검증 설계에 묶여 있으므로, 같은 프로토콜을 새 제품과 지역에 적용하고 지도를 해마다 연장할 수 있다.
